% GENERATED FILE. Edit canonical lessons, then run npm run generate:books.
% canonical-source-hash: fnv1a64:4e9b33836d9d737e
% canonical-lessons: ZH-C378-zaoyin, ZH-C378-baoyuan, ZH-C378-geyin, ZH-C378-shengren, ZH-C378-wuye

\chapter{Noise Next Door}
\label{ch:zh-noise-next-door}
\hlchaptermodality{378}

\begin{chapteropening}
\textbf{By the end of this chapter:} \emph{I can talk about noise, complaining, thin walls, strangers at the door and the management office.}

Complete the last lesson of chapter 378: I can talk about noise, complaining, thin walls, strangers at the door and the management office.
\end{chapteropening}

\section[zàoyīn]{zàoyīn — unwanted noise; din}

\label{lesson:ZH-C378-zaoyin}



\begin{quote}
Before the new one: say the Mandarin for the area around (a place), then the Mandarin for a home near a good school.
\end{quote}

\subsection*{You'll want to know: zàoyīn}
\emph{zàoyīn} — \textquotedblleft{}unwanted noise; din\textquotedblright{}.

\begin{quote}
\emph{wǎnshang wàibian de zàoyīn tài dà le} — The noise outside is far too loud at night.
\end{quote}

Its characters come later; for now it is a word for the ear.

\begin{cousinweb}
\emph{zàoyīn} is built from \emph{zào} 'clamour' and \zh{音} 'a sound', both inherited from Middle Chinese, \emph{sawH} and \emph{'im} (Guangyun, in Baxter's transcription), which glosses the first 'a crowd shouting'.
\end{cousinweb}

\begin{grammarlens}[title={What you've built}]
One more word for giving reasons.
\end{grammarlens}

\subsection*{Your turn}
\begin{itemize}
\raggedright
  \item \emph{Say it:} \emph{zàoyīn}
  \item \emph{Say it:} \emph{zàoyīn}, once more
  \item \emph{Say it:} \emph{wǎnshang wàibian de zàoyīn tài dà le} — The noise outside is far too loud at night.
  \item \emph{Recall:} say the Mandarin for the area around (a place), then the Mandarin for a home near a good school, then say \emph{zàoyīn} again
\end{itemize}

\begin{tcolorbox}[breakable,colback=teal!4,colframe=teal!35!black,title={Before you move on}]
What does zàoyīn mean? (\textquotedblleft{}Unwanted noise; din\textquotedblright{}.) Say it once more.
\end{tcolorbox}
\section[bàoyuàn]{bàoyuàn — to complain; to grumble}

\label{lesson:ZH-C378-baoyuan}



\begin{quote}
Before the new one: say the Mandarin for a home near a good school, then the Mandarin for unwanted noise.
\end{quote}

\subsection*{You'll want to know: bàoyuàn}
\emph{bàoyuàn} — \textquotedblleft{}to complain; to grumble\textquotedblright{}.

\begin{quote}
\emph{línjū bàoyuàn wǒmen tài chǎo le} — The neighbours complained that we were too noisy.
\end{quote}

Its characters come later; for now it is a word for the ear.

\begin{cousinweb}
\emph{bàoyuàn} is built from \emph{bào} 'to hold' and \emph{yuàn} 'resentment', both inherited from Middle Chinese, \emph{bawX} and \emph{'jwonH} (Guangyun, in Baxter's transcription), which glosses the first 'to grasp' and the second 'to resent': holding on to a grudge.
\end{cousinweb}

\begin{grammarlens}[title={What you've built}]
One more word for giving reasons.
\end{grammarlens}

\subsection*{Your turn}
\begin{itemize}
\raggedright
  \item \emph{Say it:} \emph{bàoyuàn}
  \item \emph{Say it:} \emph{bàoyuàn}, once more
  \item \emph{Say it:} \emph{línjū bàoyuàn wǒmen tài chǎo le} — The neighbours complained that we were too noisy.
  \item \emph{Recall:} say the Mandarin for a home near a good school, then the Mandarin for unwanted noise, then say \emph{bàoyuàn} again
\end{itemize}

\begin{tcolorbox}[breakable,colback=teal!4,colframe=teal!35!black,title={Before you move on}]
What does bàoyuàn mean? (\textquotedblleft{}To complain; to grumble\textquotedblright{}.) Say it once more.
\end{tcolorbox}
\section[géyīn]{géyīn — soundproofing; soundproof}

\label{lesson:ZH-C378-geyin}



\begin{quote}
Before the new one: say the Mandarin for unwanted noise, then the Mandarin for to grumble about.
\end{quote}

\subsection*{You'll want to know: géyīn}
\emph{géyīn} — \textquotedblleft{}soundproofing; soundproof\textquotedblright{}.

\begin{quote}
\emph{zhège fángzi de géyīn bù hǎo suǒyǐ hěn chǎo} — This house isn't soundproofed well, so it's noisy.
\end{quote}

Its characters come later; for now it is a word for the ear.

\begin{cousinweb}
\emph{géyīn} is built from \emph{gé} 'to block' and \zh{音} 'a sound', both inherited from Middle Chinese, \emph{keak} and \emph{'im} (Guangyun, in Baxter's transcription), which glosses the first 'to stop up'.
\end{cousinweb}

\begin{grammarlens}[title={What you've built}]
One more word for giving reasons.
\end{grammarlens}

\subsection*{Your turn}
\begin{itemize}
\raggedright
  \item \emph{Say it:} \emph{géyīn}
  \item \emph{Say it:} \emph{géyīn}, once more
  \item \emph{Say it:} \emph{zhège fángzi de géyīn bù hǎo suǒyǐ hěn chǎo} — This house isn't soundproofed well, so it's noisy.
  \item \emph{Recall:} say the Mandarin for unwanted noise, then the Mandarin for to grumble about, then say \emph{géyīn} again
\end{itemize}

\begin{tcolorbox}[breakable,colback=teal!4,colframe=teal!35!black,title={Before you move on}]
What does géyīn mean? (\textquotedblleft{}Soundproofing; soundproof\textquotedblright{}.) Say it once more.
\end{tcolorbox}
\section[\texorpdfstring{\zh{生人} (shēngrén)}{生人 (shēngrén)}]{\zh{生人} (shēngrén) — stranger}

\label{lesson:ZH-C378-shengren}



\begin{quote}
Before the new one: say the Mandarin for to grumble about, then the Mandarin for soundproofing (keeping sound out).
\end{quote}

\subsection*{You'll want to know: \zh{生人}}
\textbf{\zh{生人}} — \emph{shēngrén} — \textquotedblleft{}stranger\textquotedblright{}.

\begin{quote}
\textbf{\zh{晚上有生人来你最好不要开门}} \emph{wǎnshang yǒu shēngrén lái nǐ zuìhǎo bú yào kāimén} — If a stranger comes at night, you'd better not open the door.
\end{quote}

\begin{cousinweb}
\textbf{\zh{生人}} (\emph{shēngrén}) is built from \zh{生} 'raw, new' and \zh{人} 'a person', both inherited from Middle Chinese, \emph{srjaeng} and \emph{nyin} (Guangyun, in Baxter's transcription), which glosses the first 'to grow'. The sense 'unfamiliar' is not given there.
\end{cousinweb}

\begin{grammarlens}[title={What you've built}]
One more word for giving reasons.
\end{grammarlens}

\subsection*{Your turn}
\begin{itemize}
\raggedright
  \item \emph{Say it:} \emph{shēngrén}
  \item \emph{Say it:} \emph{shēngrén}, once more
  \item \emph{Say it:} \emph{wǎnshang yǒu shēngrén lái nǐ zuìhǎo bú yào kāimén} — If a stranger comes at night, you'd better not open the door.
  \item \emph{Recall:} say the Mandarin for to grumble about, then the Mandarin for soundproofing (keeping sound out), then say \emph{shēngrén} again
\end{itemize}

\begin{tcolorbox}[breakable,colback=teal!4,colframe=teal!35!black,title={Before you move on}]
What does \zh{生人} mean? (\textquotedblleft{}Stranger\textquotedblright{}.) Say it once more.
\end{tcolorbox}
\section[wùyè]{wùyè — property management; the management office}

\label{lesson:ZH-C378-wuye}



\begin{quote}
Before the new one: say the Mandarin for soundproofing (keeping sound out), then the Mandarin for a stranger.
\end{quote}

\subsection*{You'll want to know: wùyè}
\emph{wùyè} — \textquotedblleft{}property management; the management office\textquotedblright{}.

\begin{quote}
\emph{dēng huài le nǐ yīnggāi zhǎo wùyè} — If the light's broken, you should call property management.
\end{quote}

Its characters come later; for now it is a word for the ear.

\begin{cousinweb}
\emph{wùyè} is built from \emph{wù} 'a thing' and \emph{yè} 'an undertaking', both inherited from Middle Chinese, \emph{mjut} and \emph{ngjaep} (Guangyun, in Baxter's transcription), which glosses the first 'the ten thousand things' and the second 'an affair'. 'Property', and the office that looks after it, are modern senses.
\end{cousinweb}

\begin{grammarlens}[title={What you've built}]
One more word for giving reasons.
\end{grammarlens}

\subsection*{Your turn}
\begin{itemize}
\raggedright
  \item \emph{Say it:} \emph{wùyè}
  \item \emph{Say it:} \emph{wùyè}, once more
  \item \emph{Say it:} \emph{dēng huài le nǐ yīnggāi zhǎo wùyè} — If the light's broken, you should call property management.
  \item \emph{Recall:} say the Mandarin for soundproofing (keeping sound out), then the Mandarin for a stranger, then say \emph{wùyè} again
\end{itemize}

\begin{tcolorbox}[breakable,colback=teal!4,colframe=teal!35!black,title={Before you move on}]
What does wùyè mean? (\textquotedblleft{}Property management; the management office\textquotedblright{}.) Say it once more.
\end{tcolorbox}
